\documentclass[11pt]{article}
\usepackage[T1]{fontenc}
\usepackage{textcomp}
\usepackage[margin=0.75in]{geometry}
\usepackage{amsmath,amssymb,amsthm}
\usepackage{array,booktabs}
\usepackage{hyperref}
\setlength{\parindent}{0pt}
\newtheorem{theorem}{Theorem}
\theoremstyle{definition}
\newtheorem{definition}{Definition}
\title{Flow Proof Artifact: Group-inv-unique}
\author{Generated by \texttt{flow doc proof}}
\date{Flow Proof Book}
\begin{document}
\maketitle
Inverses in a group are unique.\\[0.5em]
\textbf{Source.} dummit-foote — *Abstract Algebra*, §1.1\\[0.5em]
\bigskip
\noindent{\large \textbf{Derived fact 1.} \textit{If h * g = 1 and g * k = 1 then h = k} \hfill Group \cdot inverse \cdot inverses are unique}\par
\smallskip\noindent\textit{Coordinate: Group \cdot inverse \cdot inverses are unique}\par
\smallskip\noindent\textit{Source: dummit-foote}\par
\smallskip\noindent\textit{Built on: left inverse recovers the identity, for inverse on Group, right inverse recovers the identity, for inverse on Group, parentheses do not matter, for multiplication on Group}\par
\medskip
\noindent\textbf{Goal.} If h * g = 1 and g * k = 1 then h = k\par
\noindent$\displaystyle \forall g \in Group \forall h \in Group \forall k \in Group\quad h = k$\par
\medskip
\noindent
\begin{tabular}{@{} >{\raggedright\arraybackslash}p{0.06\textwidth} >{\raggedright\arraybackslash}p{0.47\textwidth} >{\raggedleft\arraybackslash}p{0.41\textwidth} @{}}
\textbf{\#} & \textbf{Proof} & \textbf{Mathematics} \\
\midrule
\textbf{1.} & We prove that inverses are unique for inverse on Group. & \\[0.45em]
\textbf{2.} & We split into exhaustive cases — the claim must hold in each one. & \\[0.45em]
\textbf{3.} & Case 1 (see step 2): suppose h * g  equals  1. & \\[0.45em]
\textbf{4.} & Case 2 (see step 2): suppose g * k  equals  1. & \\[0.45em]
\textbf{5.} & We invoke the definitional clause governing inverse on Group: left inverse recovers the identity, for inverse on Group (instantiated for g). & \\[0.45em]
\textbf{6.} & We invoke the derived fact governing inverse on Group: right inverse recovers the identity, for inverse on Group (instantiated for g). & \\[0.45em]
\textbf{7.} & We invoke the definitional clause governing multiplication on Group: parentheses do not matter, for multiplication on Group (instantiated for h, g, k). & \\[0.45em]
\textbf{8.} & From step 4, step 5, step 6, and step 7, this implies h equals k. Together with the other cases (step 3 and step 4), the goal is discharged. Hence proven. & $\displaystyle h = k$ \\[0.45em]
\bottomrule
\end{tabular}
\label{thm:Group-inverse-inverses_are_unique}
\par\medskip\hrule
\end{document}
